\SourceLeaf{388}%
man doch einen großen Teil des Inhalts nur als Abhandlungen ansehen, in denen elementare Gegenstände eher breit als gründlich behandelt werden. Es wäre jedoch unpassend, diese Arbeiten. die einer längst vergangenen Zeit angehören, heute noch im Einzelnen zu kritisieren.

Weit gehaltvoller erscheint dagegen die neue Fortsetzung, in der Militärs, Historiker und Geografen viel Lehrreiches finden. Wir nennen hier nur die Darstellung der Schlacht bei Marengo, die Geschichte des Feldzugs in Deutschland im Jahr 1800 (die beinahe den ganzen fünften Band füllt) und die militärische Beschreibung des Einzugsgebiets der Donau, alles durch zahlreiche Karten und Pläne erläutert. Hinzu kommen die Verhandlungen einer eigens eingesetzten Kommission über die zweckmäßigste Darstellung des Geländes. Sie erörtern die Frage von vielen Seiten und veranschaulichen sie mit einer beträchtlichen Anzahl von Probezeichnungen nach unterschiedlichen Methoden. Mit Interesse wird man im Briefwechsel des Grafen \textsc{de Gisors} mit seinem Vater, dem Herzog \textsc{de Belle-Isle}, die Gespräche lesen, die der Graf ein Jahr vor Ausbruch des Siebenjährigen Krieges mit \textsc{Friedrich} dem Zweiten über militärische Fragen führte. Das gilt auch für eine Reihe bisher ungedruckter, teilweise eigenhändiger Briefe \textsc{Ludwigs} XIV., denen ein Faksimile beigegeben ist. Allerdings können sie nicht alle Leser von dem Standpunkt aus betrachten, den die Herausgeber den französischen Lesern nahelegen wollen, wenn sie in der Einleitung dazu bemerken: \textsc{Montesquieu} \textit{regarde comme le devoir de tout écrivain homme de bien de contribuer, autant qu’il est en lui, à donner à ses concitoyens des raisons d’aimer ceux à qui ils doivent obéir. Rien n’entre mieux dans cette noble pensée de} \textsc{Montesquieu}, \textit{que la publication de ces lettres et de celles qui pourront les suivre, et tout le monde reconnaîtra dans les successeurs du grand roi tout ce que son coeur avait de paternel et son âme d’héroïque.} Schließlich dürfen wir die Berichte im 3. und 4. Band über die neue große Karte von Frankreich nicht mit Stillschweigen übergehen. Ihre Ausführung wurde durch eine königliche Verordnung vom 6. August 1817 angeordnet. Zunächst wollte man die Aufnahme im Maßstab 1 zu 10000 und den Stich im Maßstab 1 zu 50000 ausführen. Dabei veranschlagte man insgesamt 611 Blätter, jedes 800 Millimeter breit und 500 Millimeter hoch, und glaubte, die ganze Arbeit in 20 Jahren vollenden zu können. Man ließ diesen Plan jedoch bald fallen und reduzierte den Maßstab für die
